\section*{Capítulo \ref{ch:b3:advanced-geochronology} --- Geocronología avanzada}

\begin{solution}{exo:b3:advanced-geochronology:1}
$t_{206} = \ln(1.04500)/\lambda_{238} = \qty{283.8}{\Ma}$; $t_{207} =
\ln(1.3300)/\lambda_{235} = \qty{289.6}{\Ma}$. Las edades difieren en
\qty{6}{\Ma}: el circón es discordante. Su $t_{207}$ es la mayor, por lo que se representa
por debajo de la \omterm{def:b3:advanced-geochronology:concordia}{concordia}, del lado de la pérdida de plomo.
\end{solution}

\begin{solution}{exo:b3:advanced-geochronology:2}
$J = (e^{5.543\times10^{-10}\times2.8201\times10^7} - 1)/1.3129 =
0.015754/1.3129 = 0.01200$. Entonces $t = \ln(1 + 0.01200\times10.000)/\lambda =
\ln 1.12/(5.543\times10^{-10}) = \qty{204.5}{\Ma}$.
\end{solution}

\begin{solution}{exo:b3:advanced-geochronology:3}
$T_{\mathrm{CHUR}} = \ln\bigl(1 + (0.51180 - 0.512638)/(0.1200 -
0.1966)\bigr)/\lambda = \ln(1.010940)/\lambda = \qty{1.67}{\Ga}$;
$T_{\mathrm{DM}} = \ln\bigl(1 + (0.51180 - 0.51315)/(0.1200 - 0.2137)\bigr)/\lambda
= \ln(1.014408)/\lambda = \qty{2.19}{\Ga}$.
\end{solution}

\begin{solution}{exo:b3:advanced-geochronology:4}
Sin desintegración, $t = N/P_0 = \num{2.0e5}/4.96 = \qty{40.3}{kyr}$. Con desintegración,
$t = -\ln(1 - N\lambda/P_0)/\lambda = -\ln(1 - 0.02014)/\lambda =
\qty{40.7}{kyr}$, un \qty{1}{\percent} más antiguo.
\end{solution}

\begin{solution}{exo:b3:advanced-geochronology:5}
Para $t$ pequeño, $x \approx \lambda_{235}t$ e $y \approx \lambda_{238}t$, por lo que
$y \approx (\lambda_{238}/\lambda_{235})x$, pendiente $0.1575$. Si el plomo se
perdió muy recientemente, $t_2 \approx 0$ y la \omterm{def:b3:advanced-geochronology:concordia}{intersección inferior} es $C(0)$, el
origen. A lo largo de una recta que pasa por el origen, $y/x$ es constante, y lo es también
$^{207}\mathrm{Pb}^*/^{206}\mathrm{Pb}^* = x/(137.8\,y)$: cada cristal proporciona
la edad plomo--plomo de la \omterm{def:b3:advanced-geochronology:concordia}{intersección superior}, la edad de formación.
\end{solution}

\begin{solution}{exo:b3:advanced-geochronology:6}
Pesos $1.5625$, $1$, $1.2346$, $2.0408$, $W = 5.838$; $\bar x =
\qty{296.09}{\Ma}$, $\sigma = 1/\sqrt W = \qty{0.41}{\Ma}$. Las desviaciones cuadráticas ponderadas
suman $2.32$; con $f = 3$, $\mathrm{MSWD} = 0.77$, por debajo de
$1 + 2\sqrt{2/3} = 2.63$: las edades concuerdan.
\end{solution}

\begin{solution}{exo:b3:advanced-geochronology:7}
$\lambda\zeta g\rho_d\rho_s/\rho_i = 1.55125\times10^{-10}\times350\times0.5
\times\num{1.10e6}\times0.5 = 0.014931$; $t = \ln(1.014931)/\lambda =
\qty{95.5}{\Ma}$. Si la roca se enfrió rápidamente, data su paso por la
zona de recocido parcial del apatito, entre unos 60 y \qty{120}{\celsius}; si
se enfrió lentamente, es una edad mixta, y las longitudes de las trazas deben aclararlo.
\end{solution}

\begin{solution}{exo:b3:advanced-geochronology:8}
Por gramo: uranio $20\times10^{-6}/238.05\times N_A = \num{5.06e16}$ átomos, de
los cuales \num{5.02e16} de uranio-238 y \num{3.64e14} de uranio-235; torio
\num{1.038e17}, $\lambda_{232} = \ln 2/\num{1.40e10} = \qty{4.95e-11}{yr^{-1}}$;
helio $10^{-9}N_A = \num{6.02e14}$. La \omterm{def:b3:advanced-geochronology:cosmogenic}{tasa de producción} es
$8\lambda_{238}{}^{238}\mathrm U + 7\lambda_{235}{}^{235}\mathrm U +
6\lambda_{232}\mathrm{Th} = (6.23 + 0.25 + 3.08)\times10^7 =
\qty{9.57e7}{\text{átomos}.g^{-1}.yr^{-1}}$, y $t = \num{6.02e14}/\num{9.57e7} =
\qty{6.29}{Myr}$. Corregida: $6.29/0.75 = \qty{8.39}{\Ma}$.
\end{solution}

\begin{solution}{exo:b3:advanced-geochronology:9}
$P_0/\lambda = 4.96/\num{4.994e-7} = \qty{9.93e6}{\text{átomos}.g^{-1}}$, alcanzada al
\qty{90}{\percent} tras $\ln 10/\lambda = \qty{4.6}{Myr}$. Con
\qty{10}{m.Myr^{-1}} $= \qty{1e-3}{cm.yr^{-1}}$: $\rho\varepsilon/\Lambda =
2.7\times10^{-3}/160 = \qty{1.69e-5}{yr^{-1}}$, $\mu = \qty{1.74e-5}{yr^{-1}}$
y $N = 4.96/\num{1.74e-5} = \qty{2.85e5}{\text{átomos}.g^{-1}}$, 35 veces menos.
\end{solution}

\begin{solution}{exo:b3:advanced-geochronology:10}
Con $W = w_1 + w_2$, $x_1 - \bar x = w_2(x_1 - x_2)/W$ y $x_2 - \bar x =
-w_1(x_1 - x_2)/W$, por lo que $\sum w_i(x_i - \bar x)^2 = (x_1 - x_2)^2(w_1w_2^2 +
w_2w_1^2)/W^2 = (x_1 - x_2)^2w_1w_2/W = (x_1 - x_2)^2/(\sigma_1^2 +
\sigma_2^2)$; con $f = 1$ esto es el MSWD. Si ambas miden la misma edad,
$x_1 - x_2$ tiene media 0 y varianza $\sigma_1^2 + \sigma_2^2$, por lo que el
valor esperado de su cuadrado dividido entre esa varianza es 1.
\end{solution}

\begin{solution}{exo:b3:advanced-geochronology:11}
$C(2700) = (13.283, 0.52019)$, $C(500) = (0.63628, 0.08065)$. $f = (0.52019 -
0.300)/(0.52019 - 0.08065) = 0.501$, la mitad de su plomo. $x = 0.499\times13.283 +
0.501\times0.63628 = 6.948$. $t_{206} = \qty{1691}{\Ma}$, $t_{207} =
\qty{2105}{\Ma}$; la razón $x/(137.8\,y) = 0.1681$ da, por iteración, una
edad plomo--plomo de \qty{2537}{\Ma}. Ninguna de las tres tiene significado: son
mezclas de los dos eventos, y solo las intersecciones las datan. La edad plomo--plomo
es un mínimo para la cristalización.
\end{solution}

\begin{solution}{exo:b3:advanced-geochronology:12}
Sin erosión: $t = -\ln(1 - \num{1.0e6}\times\num{4.994e-7}/4.96)/\lambda =
\qty{212}{kyr}$. Estado estacionario: $\rho\varepsilon/\Lambda = P_0/N - \lambda =
\qty{4.46e-6}{yr^{-1}}$, $\varepsilon = \qty{2.64e-4}{cm.yr^{-1}} =
\qty{2.6}{m.Myr^{-1}}$. Para cualquier $\varepsilon$ inferior a este valor, $N =
P_0(1 - e^{-\mu t})/\mu$ fija un $t$, desde \qty{212}{kyr} hasta
el infinito: una medida tiene dos incógnitas. Un segundo nucleido con otra
semivida, como el aluminio-26 en el mismo cuarzo, o la concentración a
profundidad en un testigo de sondeo, los separa.
\end{solution}

\begin{solution}{pb:b3:advanced-geochronology:1}
\textbf{1.} $t_{206}$ / $t_{207}$, en \unit{\Ma}: Z1 301.9 / 302.2, Z2 284.1 /
286.5, Z3 265.7 / 268.5, Z4 246.1 / 250.1, Z5 225.1 / 229.4. Un error de
\qty{0.2}{\percent} en $y$ equivale a unos \qty{0.6}{\Ma}: Z1 es concordante dentro del
error; las demás son discordantes, cada vez más discordantes desde Z2 hasta Z5.

\textbf{2.} Los circones perdieron fracciones crecientes de su plomo en un
evento posterior: se sitúan cada vez más abajo sobre la \omterm{def:b3:advanced-geochronology:concordia}{discordia}. La \omterm{def:b3:advanced-geochronology:concordia}{discordia} queda por debajo de la
\omterm{def:b3:advanced-geochronology:concordia}{concordia} cóncava, donde $x$ es mayor para un $y$ dado: $t_{207} > t_{206}$.

\textbf{3.} Z1: $X = 1/0.04794 = 20.859$, $Y = 0.3467/(137.8\times0.04794) =
0.05248$. Z5: $X = 28.145$, $Y = 0.05178$.

\textbf{4.} Pesos $1.000$, $1.235$, $1.563$, $1.563$, $2.041$ (en
$10^8$); $\bar x = 0.29402$, $\bar y = 0.040892$, $S_{xx} = \num{7.845e5}$.

\textbf{5.} $b = 0.13260$, $a = \bar y - b\bar x = 0.001905$.

\textbf{6.} A \qty{310}{\Ma}, $\mathrm dy/\mathrm dx = 0.1575\,
e^{(\lambda_{238} - \lambda_{235})t} = 0.1218$. La \omterm{def:b3:advanced-geochronology:concordia}{discordia} solo es
\qty{9}{\percent} más empinada que la \omterm{def:b3:advanced-geochronology:concordia}{concordia}: la corta con un ángulo pequeño,
y la \omterm{def:b3:advanced-geochronology:concordia}{intersección superior} será imprecisa.

\textbf{7.} Newton desde \qty{310}{\Ma}: $t = \qty{311.2}{\Ma}$.

\textbf{8.} Desde \qty{100}{\Ma}: $t = \qty{99.7}{\Ma}$.

\textbf{9.} $(y_i - a - bx_i)/\sigma_i = 0.64$, $-0.91$, $0.35$, $-0.24$,
$0.17$; $\mathrm{MSWD} = 1.435/3 = 0.48$, por debajo del límite $2.63$ para $f = 3$:
los circones se ajustan a una sola recta.

\textbf{10.} $s_b = 1/\sqrt{\num{7.845e5}} = 0.00113$.

\textbf{11.} Con $b - 2s_b = 0.13034$, la \omterm{def:b3:advanced-geochronology:concordia}{intersección superior} es
\qty{322.9}{\Ma}; con $b + 2s_b = 0.13485$, \qty{302.6}{\Ma}. La edad es
$311 \pm \qty{10}{\Ma}$ a dos desviaciones típicas.

\textbf{12.} $f = (x(311.2) - x_i)/(x(311.2) - x(99.7))$: Z1 \qty{4.7}{\percent},
Z5 \qty{41}{\percent}.

\textbf{13.} $J = 0.015754/1.3129 = 0.0119999$, $0.01200$.

\textbf{14.} Paso 1: \qty{150.0}{\Ma}; paso 4: \qty{296.6}{\Ma}; paso 10:
\qty{295.0}{\Ma}.

\textbf{15.} Pasos 4 a 10: siete pasos consecutivos, \qty{82}{\percent} del
argón-39, sin tendencia. Los pasos 1 a 3 suben en escalera.

\textbf{16.} Errores iguales: la media es $2072.9/7 = \qty{296.1}{\Ma}$, $\sigma =
1.1/\sqrt7 = \qty{0.42}{\Ma}$; la suma de los cuadrados de las desviaciones es $3.83$,
$\mathrm{MSWD} = 3.83/(1.21\times6) = 0.53$, por debajo de $2.15$.

\textbf{17.} $\bar R = \sum f_iR_i = 14.270$, $t = \qty{285.2}{\Ma}$: el
argón perdido de los bordes, liberado en los primeros pasos, reduce la media, como
lo haría con una edad potasio--argón.

\textbf{18.} $(700 - 375)/(311.2 - 296.1) = \qty{21.5}{\celsius.Myr^{-1}}$
(con un intervalo amplio, ya que la edad del circón se conoce con $\pm\qty{10}{\Ma}$).

\textbf{19.} Hacia \qty{100}{\Ma} el granito estaba enterrado bajo la
cobertera mesozoica, y un calentamiento o una circulación de fluidos hizo que los circones perdieran plomo.
La biotita perdió argón de sus bordes, lo que registra la escalera; su
paso más bajo, \qty{150}{\Ma}, es una mezcla de argón antiguo y reajustado. La meseta
se conservó: la biotita no se calentó lo suficiente para perder todo su argón.

\textbf{20.} La pendiente de mínimos cuadrados de la altitud frente a la edad es
\qty{20.1}{m.Myr^{-1}}.

\textbf{21.} $(65 - 10)/30 = \qty{1.83}{km}$; desde \qty{43.1}{\Ma}, una media de
$1830/43.1 = \qty{42}{m.Myr^{-1}}$.

\textbf{22.} $\rho\varepsilon/\Lambda = P_0/N - \lambda = \num{3.415e-5} -
\num{0.050e-5} = \qty{3.365e-5}{yr^{-1}}$; $\varepsilon = \num{3.365e-5}\times
160/2.7 = \qty{1.99e-3}{cm.yr^{-1}} = \qty{19.9}{m.Myr^{-1}}$.

\textbf{23.} $P_0\Lambda/(\rho N) = \qty{20.2}{m.Myr^{-1}}$, \qty{1.5}{\percent}
más: la desintegración apenas influye a esta tasa.

\textbf{24.} $\Lambda/(\rho\varepsilon) = 160/(2.7\times\num{1.99e-3}) =
\qty{30}{kyr}$. El granito se elevó a unos \qty{20}{m.Myr^{-1}} entre 52 y
\qty{43}{\Ma}, a \qty{42}{m.Myr^{-1}} de media desde entonces, y erosiona a
\qty{20}{m.Myr^{-1}} en la actualidad: la exhumación fue más rápida en algún momento del
Cenozoico.

\textbf{25.} El granito de Varden cristalizó hace
\textbf{\qty[text-series-to-math=true, separate-uncertainty=true]{311 \pm 10}{\Ma}} (\omterm{def:b3:advanced-geochronology:concordia}{intersección superior},
a dos desviaciones típicas).
\end{solution}
